\section{Tail Risk under Selective Labels}
\label{sec:setup}

Tail-risk certification requires a target population and loss, a criterion for large errors, and a restriction on hidden outcomes. CVaR measures average loss within a chosen upper tail, making the severity of the largest prediction errors part of the assessment. The observation model specifies what historical records identify; the sensitivity model bounds how observed and hidden outcome probabilities can differ. Together, these choices define the population risks that our certificates must cover.

\subsection{Target population and observation process}
\label{sec:target}
Let \(\mathcal F\) be the candidate predictor class and \(P\) the full-data distribution of \((X,S,Y)\), with covariate marginal \(P_X\). We evaluate \(f\in\mathcal F\) using loss \(\ell(f(X),Y)\), where \(Y\in\{0,1\}\). In decision settings, \(Y\) denotes \(Y(a^\star)\) under one specified action held fixed across the target population and predictor comparisons. The indicator \(S\in\{0,1\}\) equals one when the outcome is observed. We observe \(n\) independent copies of \(O=(X,S,Y^{\mathrm{obs}})\), where \(Y^{\mathrm{obs}}=SY\). The observed law \(P_{\mathrm{obs}}\) identifies \(P_X\) and
\[
 \pi(x)=P(S=1\mid X=x),
 \qquad p_1(x)=P(Y=1\mid X=x,S=1),
\]
while \(p_0(x)=P(Y=1\mid X=x,S=0)\) is hidden wherever \(\pi(x)<1\). Here \(\pi\) is the historical observation probability and \(p_1\) the outcome probability among observed cases; both are auxiliary to the predictor \(f\) being assessed.

Bounded loss gives a common risk scale and supports finite-sample error control. Population overlap makes the observed-outcome distribution available as a reference throughout the target covariate population. We formalize these conditions as follows.

\begin{assumption}[Bounded loss and population overlap]
\label{ass:regularity}
For every \(f\in\mathcal F\) and \(y\in\{0,1\}\), the map \(x\mapsto\ell(f(x),y)\) is measurable and takes values in \([0,1]\), and \(0<\pi_{\min}\le\pi(X)\le1\) almost surely.
\end{assumption}

Technical regularity and target-label conventions are in Appendix~\ref{app:compatibility}. Write \(z_+=\max\{z,0\}\) and \(L_f=\ell(f(X),Y)\). At level \(\alpha\in(0,1)\), \(\beta=1-\alpha\) is the upper-tail fraction. The threshold \(\eta\) is a candidate loss cutoff, and \((L_f-\eta)_+\) is the excess above it. At \(\alpha=.95\), CVaR averages the worst 5\% of population losses, allowing fractional mass at a tied cutoff:
\begin{equation}
 \rho_{\alpha,P}(f)=\operatorname{CVaR}_{\alpha,P}(L_f)
 =\inf_{\eta\in[0,1]}\left\{\eta+\beta^{-1}\mathbb E_P[(L_f-\eta)_+]\right\}.
\label{eq:cvar-definition}
\end{equation}
The expectation includes both observed and hidden cases.

\subsection{Sensitivity model and certification objective}
\label{sec:sensitivity}
Outcome-dependent observation can change label prevalence even at fixed covariates. We use the observed outcome probability as a reference and express the allowed departure of hidden outcomes through a multiplicative bound on their odds. For a finite, externally specified \(\Gamma\ge1\), impose the outcome-odds restriction \citep{scharfstein1999adjusting,tan2006distributional,zhao2019sensitivity}:
\begin{equation}
\begin{gathered}
 p_0(x)\in[\underline p_\Gamma(x),\overline p_\Gamma(x)]
 =[q_-(p_1(x)),q_+(p_1(x))],\qquad \pi(x)<1,\\
 q_-(p)=\frac{p}{p+\Gamma(1-p)},
 \qquad q_+(p)=\frac{\Gamma p}{1-p+\Gamma p}.
\end{gathered}
\label{eq:sensitivity-model}
\end{equation}
For interior probabilities, hidden-to-observed outcome odds lie between \(\Gamma^{-1}\) and \(\Gamma\); the formulas also cover \(p_1\in\{0,1\}\). At \(\Gamma=1\), \(p_0(x)=p_1(x)\): hidden and observed outcomes have the same conditional distribution given \(X\), called conditional ignorability. For example, \(p_1(x)=.5\) and \(\Gamma=2\) allow \(p_0(x)\in[1/3,2/3]\). Larger \(\Gamma\) permits greater departure. The chosen \(\Gamma\) defines a declared sensitivity scenario, supported by domain knowledge or additional outcome information (Appendix~\ref{app:gamma-calibration}).

Let \(\mathcal P_\Gamma(P_{\mathrm{obs}})\) contain all full-data laws with observed marginal \(P_{\mathrm{obs}}\) and hidden regression satisfying this restriction. Their risk endpoints are
\begin{equation}
 \underline\rho_{\alpha,\Gamma}(f)
 =\inf_{P\in\mathcal P_\Gamma}\rho_{\alpha,P}(f),
 \qquad
 \overline\rho_{\alpha,\Gamma}(f)
 =\sup_{P\in\mathcal P_\Gamma}\rho_{\alpha,P}(f).
\label{eq:population-endpoints}
\end{equation}
For fixed \(\alpha\) and \(\Gamma\), our statistical objective is an interval covering every compatible risk simultaneously for every \(f\in\mathcal F\), with probability at least \(1-\delta_{\mathrm{tot}}\) over the observed sample. Here \(\delta_{\mathrm{tot}}\in(0,1)\) is the total failure budget; Section~\ref{sec:learning} accounts for sample fluctuation and auxiliary-estimation error within this budget. Thus \(\Gamma\) specifies the admissible hidden outcomes, while \(\delta_{\mathrm{tot}}\) specifies the statistical guarantee.

We next derive these population risk endpoints from the observed law and show that every value between them is attainable.
